% !TEX root = ../main.tex
\section{Introduction}

Linear programming (LP) has been the workhorse of strategic forest planning for half a century. Since the Model~I and Model~II formulations of \citet{johnson1977techniques}, LP harvest-scheduling models, such as the USDA Forest Service's FORPLAN system and its successors, have been used to project long-term timber supply, to test harvest-flow policies, and to support trade-off analysis through the shadow prices that an optimal LP solution provides. In nearly all of such applications, the model is solved once, over the full multi-decade planning horizon, as an \emph{open-loop} problem: the entire sequence of future decisions is chosen at time zero, and the resulting schedule of harvest levels and residual forest structure is presented as the plan.

Real-world forest planning likely does not align with the open-loop assumption that today's planner could precommit to every future decision. Real planning is sequential: each period, the planner observes the realized forest state and re-solves the planning problem from that state, under the same objectives and constraints. If the tail of the original plan is not optimal for the subproblem that starts at the realized state, the re-solving planner will deviate from it. The announced plan then fails Bellman's principle of optimality \citep{bellman1957dynamic} and is \emph{dynamically inconsistent}. The projected harvest trajectory, residual growing stock, and shadow prices are not credible statements about what will or should happen. A plan that will not be followed cannot be a sound basis for long-term sustainable policies. Comparative analyses performed on it are biased, including shadow-price trade-off analyses that often motivate such modeling studies.

Dynamic inconsistency is also a problem between generations. An open-loop plan can raise today's harvest by scheduling harvests that later planners will no longer find worth taking \citep[p.~332]{gunn2007models}: inconsistent decisions are selected for a value that the first planner captures and later planners treat as sunk---an ``income transfer from future planners to the current planner'' \citep[pp.~41, 142]{daugherty1991credibility}. Intergenerational equity in LP-based forest planning has also been addressed directly by \citet{church1999considering}. Dynamic inconsistency is a general phenomenon in dynamic economics \citep{strotz1956myopia,kydland1977rules}, and it was demonstrated for forest planning more than three decades ago: in his 1991 PhD thesis at the University of California, Berkeley, P. J. Daugherty proved analytically and demonstrated numerically that LP-based forest-planning models solved as open-loop formulations admit dynamically inconsistent plans, identified the structural conditions under which the inconsistency arises, and characterized its behavior over a wide range of initial forest conditions and harvest policies \citep{daugherty1991credibility}. The thesis remains, to our knowledge, the most complete treatment of the problem in the forest-planning literature; yet its analysis was never published as a journal article. The phenomenon itself has surfaced in the literature under various names. \citet{mcquillan1986declining} named the ``declining even-flow effect'' in national forest planning; \citet[pp.~331--332]{gunn2007models}, citing \citet{daugherty1991credibility}, calls it the ``declining non-declining yield''; and \citet{paradis2013risk} documented the systematic drift that arises between planned and implemented harvest under incoherent hierarchical planning, citing \citet{daugherty1991credibility} directly and using ``credibility'' in the thesis's sense. Yet the original analysis remains largely unread, and the field has lacked an open, citable, reproducible reference implementation against which to check it.

This paper formalizes dynamic inconsistency for LP forest plans and provides that reference, an open reproduction and refinement of the methods and case study of \citet{daugherty1991credibility} built on the open-source \texttt{ws3} wood-supply model \citep{paradis2026ws3} and the public-domain Umpqua planning documents. The contribution is fourfold. First, we give an operational divergence metric and occurrence criterion for dynamic inconsistency, making the phenomenon measurable and comparable across models. Second, we provide an open reproduction package built on \texttt{ws3}: the case-study data, the LP formulations, the sequential-replanning simulator, and the scenario grid. It includes a new \texttt{ws3} constraint form for the classic FORPLAN harvest-flow policies (non-declining yield and bounded sequential flow), released in \texttt{ws3} v1.1.0a5. Third, we reproduce and refine the thesis's empirical characterization: where the thesis ran 178 selected combinations of its factors, we run the full factorial of its 18 initial forest conditions, 4 discount rates and 6 harvest-flow policies (432 scenarios), cleanly separating the harvest-flow constraint's role against a flow-unconstrained control and showing that inconsistency occurs at every discount rate while its magnitude depends on the rate. Fourth---moving beyond reproduction---we run four extension experiments with alternative formulations, each asking the same question: \textit{Does this formulation mitigate, or eliminate, the dynamic inconsistency?} The experiments assess the alternative \emph{discount-rate shapes} (declining-rate schedules), the \emph{flow constraint's denomination} (net revenue instead of volume),  the \emph{flow constraint's shape} (rolling-mean rather than point-wise linkages), and a \emph{maximum flow-calibration} alternative that replaces the even-flow constraints with an iteratively tightened max-harvest constraint. We emphasize the framing: the underlying science is that of \citet{daugherty1991credibility}; the novelty is the operational metric, the open and reproducible benchmark, the rigorous driver diagnosis, and the systematic mitigation tests that the open reproduction package makes possible.
